\section{Lorentz Invariance (P3, P11, P27)}
\label{sec:lorentz}

\subsection{The both-worlds substrate (P3)}

The substrate must be two things that pull against each other: Lorentz-safe (no
preferred frame, which random structures give but rigid ones do not) and navigable
(routable from local information, which addressable rigid structures give but generic
random ones do not). \textbf{Where we landed:} the connected hyperbolic random graph is
both at once. With mean degree about eleven it has exponential reach, anisotropy
$0.07$ (as Lorentz-safe as a pure sprinkling), and $100$ percent backtracking
navigability, and all three survive an eightfold growth of the mesh at constant
density. The price is greedy geometric routing instead of exact addressing, consistent
with no hidden state. This is the result that fixes the chosen substrate.

\subsection{The dynamics is Lorentz-isotropic in the infrared (P11)}

Does the emergent dynamics propagate isotropically? \textbf{Where we landed, with an
honest correction:} both a lattice and a random mesh propagate isotropically at long
wavelength. This is the lattice-field-theory fact that rotational and Lorentz invariance
emerge in the infrared even on a lattice, the anisotropy being a lattice-scale effect.
So the clean substrate-level distinction is P3, not the low-energy dynamics, and
discreteness threatens Lorentz invariance only at the Planck scale.

\subsection{Random discreteness is Lorentz-safe, a lattice is not (P27)}

At the discreteness scale the difference is decisive. A lattice has energy-dependent,
directional Lorentz violation: its group-speed anisotropy grows with energy (from
$0.003$ to $0.71$ across the band) and the speed of light falls below one. A random
sprinkling has none: its nearest-neighbour link directions are isotropic (a $4$-fold
anisotropy of $0.04$ versus $1.0$ for the lattice). \textbf{Where we landed:} because
the framework uses the random sprinkling, it predicts \emph{no} first-order Lorentz
violation, which is exactly what gamma-ray-burst photon timing observes
(Section~\ref{sec:predictions}). Random discreteness preserves Lorentz invariance where
regular discreteness breaks it, and the residual effect of discreteness is not coherent
violation but the stochastic swerve of Section~\ref{sec:predictions}.
